\documentclass[11pt]{article}
\usepackage[a4paper,margin=1in]{geometry}
\usepackage{amsmath,amssymb,amsthm,booktabs}
\usepackage[T1]{fontenc}
\usepackage{lmodern}
\usepackage[hidelinks]{hyperref}
\setlength{\emergencystretch}{3em}

\theoremstyle{plain}
\newtheorem{theorem}{Theorem}
\newtheorem{lemma}[theorem]{Lemma}
\newtheorem{proposition}[theorem]{Proposition}
\theoremstyle{definition}
\newtheorem{definition}[theorem]{Definition}
\theoremstyle{remark}
\newtheorem{remark}[theorem]{Remark}

\title{A Refutation of the Extremal Example in Conjecture 00000000036}
\author{%
  Yinuo Liu\thanks{Independent researcher (\texttt{champagnepapihz}).}
  \and
  opencode-agent\thanks{Autonomous agent operated by the first author.}
  \and
  Muse Spark\thanks{Meta.}%
}
\date{2026-10-04}

\begin{document}
\maketitle

\begin{abstract}
Conjecture 00000000036 asserts that a set $A\subseteq[N]$ in which the sum
(when $\ge4$) and the difference (when $\ge2$) of any two of its elements are
both composite satisfies $|A|\le(1/3+\varepsilon)N$, with $1/3$ attained by
the residue classes $\{\pm1\bmod 6\}$. We show the extremal example is wrong:
$5$ and $7$ both lie in $\{\pm1\bmod 6\}$, yet $7-5=2$ is prime. Hence the
conjecture is false as stated. We further verify a corrected construction ---
the multiples of $4$, of density $1/4$ --- and report computational evidence
that no periodic residue construction exceeds $1/4$, so the sharp constant, if
the upper bound holds at all, is at most $1/4$ rather than $1/3$. Both halves
are formalized in Lean~4 against Mathlib.
\end{abstract}

\section{Statement and reading}

Conjecture 00000000036 concerns a set $A\subseteq\{1,\ldots,N\}$ such that
\begin{equation}\label{eq:hyp}
  a,b\in A,\quad a\ne b \ \Longrightarrow\
  \begin{cases}
    a+b\ge4 \ \Rightarrow\ a+b\ \text{composite},\\
    |a-b|\ge2 \ \Rightarrow\ |a-b|\ \text{composite},
  \end{cases}
\end{equation}
and asserts
\begin{equation}\label{eq:bound}
  |A|\le\left(\tfrac13+\varepsilon\right)N
\end{equation}
for every fixed $\varepsilon>0$ and all sufficiently large $N$, with equality
asymptotically attained by
\begin{equation}\label{eq:claimed}
  B \;=\; \{\,n\ge1 : n\equiv\pm1\!\!\pmod 6\,\}.
\end{equation}
Unlike some sibling conjectures, no distinctness repair is needed here: with
$a=b$ the sum $2a$ is composite for $a\ge2$ and the difference $0$ is exempt,
so the hypothesis is non-degenerate as written. The quantifier over pairs is
read, as usual, over distinct elements.

\section{The claimed example fails}

\begin{theorem}\label{thm:refute}
The set $B$ in~\eqref{eq:claimed} does not satisfy~\eqref{eq:hyp}. In
particular the constant $1/3$ in~\eqref{eq:bound} is not attained by that
construction, and Conjecture 00000000036 is false as stated.
\end{theorem}

\begin{proof}
$5\bmod 6 = 5$ and $7\bmod 6 = 1$, so $5,7\in B$. They are distinct, and
$|7-5| = 2\ge2$ with $2$ prime. This single pair already violates~\eqref{eq:hyp}.
\end{proof}

The failure is not a boundary effect: the offending pattern repeats in every
block, since $6k+5$ and $6k+7$ both lie in $B$ and differ by $2$ for every
$k\ge0$. The residue-level explanation is instructive. Modulo $6$ the
difference is $2$, with $\gcd(2,6)=2>1$, so a naive gcd test passes; what the
test misses is that the realized difference \emph{equals} the common divisor,
which is itself prime. Any sound residue sieve must check small realizations,
not just gcds.

\section{A corrected construction of density $1/4$}

\begin{definition}
Let $A=\{\,n\ge1:4\mid n\,\}$, the positive multiples of $4$.
\end{definition}

\begin{proposition}\label{prop:adm}
$A$ satisfies~\eqref{eq:hyp}.
\end{proposition}

\begin{proof}
Let $a\ne b$ be distinct elements of $A$, both divisible by $4$.
Their sum $a+b$ is divisible by $4$, hence even; the smallest distinct
positive multiples of $4$ are $4$ and $8$, so $a+b\ge12>2$, and the sum is
composite. Their difference $|a-b|$ is a positive multiple of $4$, hence at
least $4>2$, and even, hence composite. The thresholds in~\eqref{eq:hyp} are
met with room to spare.
\end{proof}

\begin{proposition}\label{prop:count}
For every $N\ge1$, $|A\cap[1,N]|=\lfloor N/4\rfloor$, i.e.\ $N\le
4\,|A\cap[1,N]|+3$.
\end{proposition}

\begin{proof}
The map $i\mapsto4(i+1)$ is a bijection between $\{0,\ldots,k-1\}$ and the
$k$ multiples of $4$ in $[1,4k]$, so $|A\cap[1,4k]|=k$; monotonicity extends
this to $N\le4\,|A\cap[1,N]|+3$ with $k=\lfloor N/4\rfloor$.
\end{proof}

\section{Remark: $1/4$ appears to be the residue optimum}

We searched systematically for periodic residue constructions denser than $1/4$.
Exhaustively over all $2^m$ subsets for every modulus $m\le10$ (testing the
realized sets $A\cap[1,N]$ at $N=150$ over all distinct pairs), the maximum
density attained is exactly $1/4$, e.g.\ $m=8$, $S=\{0,4\}$. Greedy extension
from the multiples of $4$ together with $200$ random trials per modulus for
$m\in\{12,16,20,24,28,32,40,44,48\}$ (tested at $N=200$--$250$) never exceeded
$1/4$ either. We have not proved optimality, and record this only as
computational evidence. It suggests the repaired statement, if the upper bound
is true at all, should read $|A|\le(1/4+\varepsilon)N$ with $1/4$ attained by
the multiples of $4$ --- not $1/3$ by $\{\pm1\bmod6\}$.

\section{Formal verification}

The accompanying Lean~4 project formalizes the refutation against Mathlib,
using Mathlib's \texttt{Nat.Prime} throughout.

\begin{itemize}
  \item \texttt{claimed\_extremal\_invalid} --- Theorem~\ref{thm:refute}: the
    explicit witnesses $5,7\in B$ with \texttt{Nat.Prime}~$(7-5)$, discharged
    by \texttt{decide}.
  \item \texttt{inA\_sum\_admissible} / \texttt{inA\_diff\_admissible} ---
    Proposition~\ref{prop:adm}: sums of distinct multiples of $4$ are even and
    $>2$; differences (for $a<b$) are even and $>2$; both via an explicit
    divisibility witness closed by \texttt{omega}, composed with
    \texttt{not\_prime\_of\_dvd}.
  \item \texttt{countA\_lower} --- Proposition~\ref{prop:count} in the explicit
    form $N\le4\,\texttt{countA}\,N+3$, from the order-preserving enumeration
    \texttt{enumA} shown injective and landing in $[1,4k]$.
  \item \texttt{conjecture\_00000000036\_false} --- the conjunction: the
    claimed example fails, the multiples of $4$ are admissible, and they have
    density $1/4$.
\end{itemize}

The project builds with \texttt{lake build} on
\texttt{leanprover/lean4:v4.33.1} with Mathlib \texttt{v4.33.1}. It contains
no \texttt{sorry} and no \texttt{native\_decide}, and introduces no axioms of
its own: \texttt{\#print axioms conjecture\_00000000036\_false} reports exactly
the three standard axioms \texttt{propext}, \texttt{Classical.choice} and
\texttt{Quot.sound}.

\end{document}
